\section{American Response: Mercury, Gemini, Apollo}

\subsection{NASA and the crash program}
Faced with repeated Soviet firsts, the United States organized an unprecedented peacetime mobilization of engineering talent. NASA grew rapidly, contracting with universities and thousands of companies to build launch vehicles, spacecraft, tracking networks, and test facilities almost simultaneously. The result was a crash program whose scale rivaled the largest industrial efforts in the nation's history.

Project Mercury came first, proving that Americans could survive launch, orbit, and reentry. Its successes were hard-won and public, and they restored a measure of confidence even as they revealed how much remained unknown about living and working in space.

Managing such a program demanded new methods. Systems engineering, rigorous testing schedules, and tight configuration control became as important as any single rocket, because coordinating so many contractors under deadline pressure was itself a formidable technical challenge. \cite{logsdon2010} \cite{nasa2019}

\subsection{From Earth orbit to the lunar landing goal}
In May 1961 President Kennedy committed the United States to landing a man on the Moon and returning him safely before the decade's end. The pledge was audacious: the country had barely fifteen minutes of human spaceflight experience, yet it now had a deadline and a destination that the whole world could understand.

Choosing the Moon was a strategic decision as much as a scientific one. It was far enough ahead that the Soviet lead in early milestones would not decide the outcome, and dramatic enough that success would be unmistakable. The goal concentrated national effort on a single, measurable objective.

Project Gemini bridged the gap between ambition and capability. Its flights practiced the rendezvous, docking, long-duration flight, and spacewalking that a lunar mission would require, turning Kennedy's promise into a sequence of solvable engineering problems. \cite{logsdon2010} \cite{cadbury2006}

The ideal rocket equation relates velocity change to mass ratio:
\begin{equation}
\Delta v = v_e \ln\!\left(\frac{m_0}{m_f}\right)
\label{eq:rocket}
\end{equation}
where $v_e$ is effective exhaust velocity, $m_0$ initial mass, and $m_f$ final mass.

\begin{figure}[htbp]
\centering
\includegraphics[width=0.88\linewidth]{figures/ch03_apollo.jpg}
\caption{Buzz Aldrin on the Moon during Apollo 11.}
\label{fig:ch_buzz_aldrin_on_the_moon_during_a}
\end{figure}

\bigskip
\begin{otherlanguage}{hebrew}
\subsection*{סיכום הפרק}

פרק זה מתאר את תוכניות החלל האמריקאיות - מרקורי, ג'מיני ואפולו - שריכזו משאבים לאומיים כדי להשיג נחיתה על הירח לפני ברית המועצות.

\end{otherlanguage}

\clearpage
